\documentclass[11pt]{amsart}

\usepackage[T1]{fontenc}
\usepackage{lmodern}
\usepackage{microtype}
\usepackage{amsmath,amssymb}
\usepackage[hidelinks]{hyperref}
\hypersetup{
  pdftitle={A Counterexample to Samart's Half-Shift Formula for n4},
  pdfsubject={Mahler measures and Eisenstein--Kronecker series},
  pdfkeywords={Mahler measure, Eisenstein--Kronecker series, CM point, validated numerics}
}

\newtheorem{theorem}{Theorem}
\newtheorem{lemma}[theorem]{Lemma}

\newcommand{\Z}{\mathbb Z}
\newcommand{\e}{\mathrm e}
\newcommand{\EK}{\mathcal E_4}
\DeclareMathOperator{\arcosh}{arcosh}
\renewcommand{\Re}{\operatorname{Re}}
\renewcommand{\Im}{\operatorname{Im}}

\title[A counterexample to Samart's half-shift formula]
      {A Counterexample to Samart's Half-Shift Formula for $n_4$}

\date{}

\subjclass[2020]{Primary 11F67; Secondary 11Y60}
\keywords{Mahler measure, Eisenstein--Kronecker series, CM point,
validated numerics}

\begin{document}

\begin{abstract}
Proposition~3.1(iii) of Samart's paper asserts an Eisenstein--Kronecker
formula for $n_4(s_4(\tau))$ when
$\tau=\frac12+iy$ and $y>\frac12$. At
$\tau_0=\frac12+\frac{\sqrt6}{4}i$ we prove, by interval arithmetic,
that
\[
 n_4\bigl(s_4(\tau_0)\bigr)
 \ge \frac{151680651}{67108864}>2.26,
 \qquad
 \mathcal E_4(\tau_0)<2.22516,
\]
where $\mathcal E_4$ is the lattice sum in the proposition. Thus the
half-shift clause of Proposition~3.1(iii) is false. Since
$\Im\tau_0=\sqrt6/4<1/\sqrt2$, it is that clause alone that is
contradicted; the half-shift range was introduced in the 2016 form of
the proposition. The proof uses an
exact two-dimensional Jensen reduction and a finite interval
certificate.
\end{abstract}

\maketitle

\section{The counterexample}

For a Laurent polynomial $P$ in $r$ variables, let
\[
 m(P)=\int_{[0,1]^r}
 \log\left|P\bigl(\e^{2\pi i t_1},\ldots,\e^{2\pi i t_r}\bigr)\right|
 \,dt_1\cdots dt_r
\]
and define
\[
 n_4(s)=4m\bigl(x^4+y^4+z^4+1+s^{1/4}xyz\bigr).
\]
The value is independent of the fourth root, by a unit-modulus change
of one torus variable.

Write
\[
 \eta(\tau)=\e^{\pi i\tau/12}
 \prod_{n\ge1}\bigl(1-\e^{2\pi i n\tau}\bigr),
\]
and set
\begin{align*}
 A(\tau)&=\frac{\eta(\tau)\eta(4\tau)^2}{\eta(2\tau)^3},\\
 c(\tau)&=\left(\frac{\eta(2\tau)}{\eta(\tau)}\right)^6
 \left(16A(\tau)^4+A(\tau)^{-4}\right),
 &s_4(\tau)&=c(\tau)^4.
\end{align*}
This is Samart's parameter $s_4(q)$ with $q=\e^{2\pi i\tau}$.
For $z\ne0$, put
\[
 F(z)=\frac{4(\Re z)^2}{|z|^6}-\frac1{|z|^4},
 \qquad
 T_d(\tau)=\sum_{(m,n)\in\Z^2}^{\prime}F(dm\tau+n),
\]
where the prime omits $(0,0)$, and define
\[
 \EK(\tau)=\frac{10\Im\tau}{\pi^3}
 \bigl(-T_1(\tau)+4T_2(\tau)\bigr).
\]
The sums converge absolutely. Samart's Proposition~3.1(iii) states
that
\begin{equation}\label{eq:samart}
 n_4\bigl(s_4(\tau)\bigr)=\EK(\tau)
\end{equation}
when either $\tau=iy$ with $y\ge1/\sqrt2$, or
$\tau=\frac12+iy$ with $y>\frac12$%
~\cite[Proposition~3.1(iii)]{Samart2016}.

\begin{theorem}\label{thm:main}
Let
\[
 \tau_0=\frac{2+\sqrt{-6}}4
       =\frac12+\frac{\sqrt6}{4}i.
\]
Then
\[
 n_4\bigl(s_4(\tau_0)\bigr)
 \ge \frac{151680651}{67108864}>2.26,
 \qquad
 \EK(\tau_0)<2.22516.
\]
Thus \eqref{eq:samart} fails at a point in the half-shift range of
Proposition~3.1(iii).
\end{theorem}

Although $\tau_0$ also appears in the conjectural Table~6 of
\cite{Samart2015}, the theorem concerns only Proposition~3.1(iii).
The half-shift clause is not idle in \cite{Samart2016}, however: the
proof of \cite[Thm.~1.2(iii)]{Samart2016} in the case $s<0$ applies it
at $(\tau-1)/2\tau$ with $\tau=iy$, a point whose real part is
$\frac12$ and whose imaginary part is $1/2y>\frac12$, and $\tau_0$ is
of exactly this form, with $y=2/\sqrt6$. We do not pursue the
consequences for that theorem, or for Table~6, here.

\medskip
\noindent
\textbf{Relation to earlier work.}
That \eqref{eq:samart} fails below $\Im\tau=1/\sqrt2$ is not new.
Zheng \cite[\S5]{Zheng2026} shows that the lattice sum leaves the
geometric sheet there and states that below that threshold ``the
identity fails everywhere'', on numerical evidence along the lines
$\Re\tau=0$ and $\Re\tau=1/4$, and also reports high-precision
numerical evidence that the Table~6 identity $n_4(81)=40L'(g_7,0)$
fails. Since $\Im\tau_0=\sqrt6/4<1/\sqrt2$, that blanket statement
already covers $\tau_0$, so what is new here is not the phenomenon but
the target and the mode of proof. Zheng quotes the identity only in its
earlier form \cite[Prop.~2.1(iii)]{Samart2013}, for
$\Im\tau\ge1/\sqrt2$, and records that below that threshold the
identity ``is not asserted''. The point of the present note is that it
\emph{is} asserted: \cite[Prop.~3.1(iii)]{Samart2016}, which
\cite{Samart2016} introduces as a ``modified version'' of
\cite[Prop.~2.1]{Samart2013} and justifies by the remark that the
admitted values of $\tau$ ``lie in a corresponding fundamental domain
for $\Gamma_0(N)^*$'', covers $\tau=\frac12+iy$ for every $y>\frac12$,
and hence the whole strip $\frac12<y<\frac1{\sqrt2}$.
Theorem~\ref{thm:main} is thus a proof, rather than numerical evidence,
that the formula fails at a point of the range the proposition does
assert; membership of such a fundamental domain is not by itself
sufficient for it.

\section{Jensen reduction}

\begin{lemma}\label{lem:jensen}
For $c\in\mathbb C$, set
\[
 a_c(\theta,\phi)
 =\e^{i(\theta+\phi)/2}\bigl(\e^{i\theta}+\e^{i\phi}+c\bigr),
 \qquad
 J(a)=\arcosh\left(\frac{|a-2|+|a+2|}{4}\right),
\]
where $\arcosh$ is the nonnegative real inverse of $\cosh$. Then
\begin{equation}\label{eq:jensen}
 n_4(c^4)=\frac4{(2\pi)^2}
 \int_0^{2\pi}\!\int_0^{2\pi}
 J\bigl(a_c(\theta,\phi)\bigr)\,d\theta\,d\phi.
\end{equation}
\end{lemma}

\begin{proof}
The argument of $\arcosh$ is at least $1$ by the triangle inequality,
and $J(-a)=J(a)$. Divide
$x^4+y^4+z^4+1+cxyz$ by $xyz$ and make the torus substitution
\[
 u=\frac{x^3}{yz},\qquad
 v=\frac{y^3}{xz},\qquad
 w=\frac{z^3}{xy}.
\]
Its exponent matrix has determinant $16$, so the map is a finite
torus covering and preserves normalized Haar measure. Since
$uvw=xyz$, the transformed Laurent polynomial is
\[
 u+v+w+(uvw)^{-1}+c.
\]
For fixed $u=\e^{i\theta}$ and $v=\e^{i\phi}$, Jensen's formula in $w$
reduces the integral to the roots of
\[
 w^2+(u+v+c)w+(uv)^{-1}.
\]
After a unit-modulus rescaling this is
$z^2+a_c(\theta,\phi)z+1$; the choice of square root changes $a_c$ to
$-a_c$ and hence does not change $J$.

Let $\rho$ be a root with $|\rho|\ge1$. Jensen's contribution is
$t=\log|\rho|$. Since $-a_c=\rho+\rho^{-1}$,
\[
 |a_c-2|+|a_c+2|
 =\frac{|\rho+1|^2+|\rho-1|^2}{|\rho|}
 =4\cosh t.
\]
Thus the contribution equals $J(a_c)$. The logarithm of a nonzero
Laurent polynomial is integrable on the torus, so Fubini's theorem and
Jensen's formula apply, proving \eqref{eq:jensen}.
\end{proof}

\section{Validated bounds}

Each inequality below is checked by a finite interval computation from
the constants printed here.

\subsection*{The eta parameter}

Let $q=\e^{2\pi i\tau}$, $r=|q|$, and
$R_N=\prod_{n>N}(1-q^n)$. Then
\[
 |\log R_N|
 \le \sum_{n>N}\frac{r^n}{1-r^n}
 \le \frac{r^{N+1}}{(1-r)(1-r^{N+1})},
\]
and hence
\begin{equation}\label{eq:eta-tail}
 |R_N-1|
 \le \exp\left(\frac{r^{N+1}}
 {(1-r)(1-r^{N+1})}\right)-1.
\end{equation}
Using $N=30$ in the eta products at
$\tau_0,2\tau_0,4\tau_0$, with the defining factor
$\e^{\pi i\tau/12}$ evaluated directly from $\tau$, gives
\begin{equation}\label{eq:c-box}
 0.88<\Re c(\tau_0)<0.89,
 \qquad
 -0.89<\Im c(\tau_0)<-0.88.
\end{equation}
Here $\Re c(\tau_0)=0.8850840762662\ldots$ and
$\Im c(\tau_0)=-0.8850840762662\ldots$.

\subsection*{The lattice sum}

The partial-fraction identity
\[
 \sum_{n\in\Z}(z+n)^{-3}
 =\frac{\pi^3\cos(\pi z)}{\sin^3(\pi z)}
\]
gives, after pairing the nonzero lattice rows,
\begin{align*}
 U_j(\tau)&=2\pi^3\sum_{m\ge1}
 \frac{\cos(j\pi m\tau)}{m\sin^3(j\pi m\tau)},\\
 \EK(\tau)&=\Im\left(2\pi\tau+
 \frac{10}{\pi^3}\bigl(U_1(\tau)-2U_2(\tau)\bigr)\right).
\end{align*}
Indeed,
\[
 \Im\left(\frac{(jm\tau+n)^{-3}}m\right)
 =-j\Im(\tau)F(jm\tau+n),
\]
and the omitted $m=0$ row of each $T_j$ is
$6\zeta(4)=\pi^4/15$.

Writing $y=\Im\tau_0$ and $q_j=\e^{-2j\pi y}$, the tail after $M$
terms satisfies
\begin{equation}\label{eq:u-tail}
 |U_j-U_{j,M}|
 \le
 \frac{8\pi^3(1+q_j^{M+1})q_j^{M+1}}
 {(M+1)(1-q_j)(1-q_j^{M+1})^3}.
\end{equation}
Interval evaluation of the first $20$ terms, together with
\eqref{eq:u-tail}, yields
\begin{equation}\label{eq:ek-bound}
 \EK(\tau_0)<2.22516.
\end{equation}
The enclosure obtained is $\EK(\tau_0)=2.2251552127\ldots$, the tail
\eqref{eq:u-tail} at $M=20$ being below $10^{-34}$.

\subsection*{The Mahler measure}

Let
\[
 c_*=\frac{177}{200}(1-i),
 \qquad
 \delta=\frac{123}{3125}.
\]
From \eqref{eq:c-box},
\[
 |c(\tau_0)-c_*|<\frac{\sqrt2}{200},
 \qquad
 |c(\tau_0)|<\frac{63}{50}.
\]
Partition $[0,2\pi]^2$ into $512^2$ equal squares with centers
\[
 \theta_j=\frac{2\pi(j+1/2)}{512},
 \qquad 0\le j<512.
\]
Since
\[
 |\partial_\theta a_c|,|\partial_\phi a_c|
 \le 2+\frac{|c|}{2}\le\frac{263}{100},
\]
every value in the $(j,k)$-cell lies within
\[
 \frac{263\pi}{25600}+\frac{\sqrt2}{200}<\delta
\]
of $a_{c_*}(\theta_j,\theta_k)$. The last inequality follows, for
example, from $\pi<355/113$ and $\sqrt2<99/70$.

Directed interval arithmetic gives lower bounds $\ell_{jk}^{\pm}$ for
$|a_{c_*}(\theta_j,\theta_k)\pm2|$. With $x_+=\max\{x,0\}$, define
\[
 X_{jk}=\max\left\{1,
 \frac{(\ell_{jk}^{-}-\delta)_+
       +(\ell_{jk}^{+}-\delta)_+}{4}\right\},
\]
and let $r_{jk}$ be the largest nonnegative integer such that
\[
 \cosh\left(\frac{r_{jk}}{4096}\right)\le X_{jk}.
\]
A direct interval computation over the $512^2$ cells gives the exact
sum
\begin{equation}\label{eq:sum}
 \sum_{j,k=0}^{511}r_{jk}=606722604.
\end{equation}
Of the $262144$ cells, $31988$ have $X_{jk}=1$ and hence $r_{jk}=0$;
the other $230156$ satisfy $8\le r_{jk}\le5121$, with $r_{0,0}=4050$,
$r_{255,255}=1803$, $r_{511,511}=4080$ and $\sum_kr_{0,k}=1687754$. No
$4096\arcosh(X_{jk})$ with $r_{jk}>0$ lies within $10^{-6}$ of an
integer, so the floors are insensitive to the binary64 rounding.
By Lemma~\ref{lem:jensen} and monotonicity of $\arcosh$,
\[
 n_4\bigl(s_4(\tau_0)\bigr)
 \ge \frac{4}{512^2\cdot4096}\sum_{j,k}r_{jk}
 =\frac{151680651}{67108864}>2.26.
\]
Together with \eqref{eq:ek-bound}, this proves
Theorem~\ref{thm:main}.

\medskip
\noindent
\textbf{Reproducibility.}
The recorded run used Python~3.13.5, NumPy~2.3.5, and mpmath~1.3.0.
All transcendental quantities were enclosed with the mpmath interval
context at $80$ decimal digits. The subsequent binary64 stage assumes
IEEE~754 round-to-nearest arithmetic and correctly rounded square root;
each result is enlarged outward with \texttt{numpy.nextafter}. These
semantics form the trusted computing base. Once the cell inequalities
are checked, \eqref{eq:sum} gives the Mahler lower bound as an exact
rational number.

\begin{thebibliography}{9}

\bibitem{Samart2013}
D.~Samart,
\emph{Three-variable Mahler measures and special values of modular and
Dirichlet $L$-series},
Ramanujan J. \textbf{32} (2013), no.~2, 245--268.
arXiv:1205.4803.

\bibitem{Samart2015}
D.~Samart,
\emph{Mahler measures as linear combinations of $L$-values of multiple
modular forms},
Canad. J. Math. \textbf{67} (2015), no.~2, 424--449.
\href{https://doi.org/10.4153/CJM-2014-012-8}
{doi:10.4153/CJM-2014-012-8}.

\bibitem{Samart2016}
D.~Samart,
\emph{The elliptic trilogarithm and Mahler measures of $K3$ surfaces},
Forum Math. \textbf{28} (2016), no.~3, 405--423.
\href{https://doi.org/10.1515/forum-2014-0045}
{doi:10.1515/forum-2014-0045}.
arXiv:1309.7730v2, whose numbering is followed here.

\bibitem{Zheng2026}
H.~Zheng,
\emph{Samart's conjecture $n_4(81)=40M_7$: the exact CM evaluation and
the two obstructions. A status report},
arXiv:2608.02265v1 [math.NT], 2026.

\end{thebibliography}

\end{document}